% !TEX program = pdflatex
\documentclass[11pt]{article}

\usepackage[a4paper,margin=2.5cm]{geometry}
\usepackage{amsmath,amssymb,amsthm,mathtools,bm}
\usepackage{hyperref}
\usepackage{microtype}
\usepackage{enumitem}
\usepackage{xcolor}
\usepackage{tcolorbox}
\usepackage{booktabs}
\usepackage{array}
\tcbuselibrary{breakable,skins}

% ---------- Notation ----------
\newcommand{\R}{\mathbb{R}}
\newcommand{\E}{\mathbb{E}}
\newcommand{\PP}{\mathbb{P}}
\newcommand{\cN}{\mathcal{N}}
\newcommand{\grad}{\nabla}
\newcommand{\diver}{\mathrm{div}}
\newcommand{\tr}{\mathrm{tr}}
\newcommand{\Id}{\mathrm{Id}}
\newcommand{\norm}[1]{\left\lVert #1 \right\rVert}
\newcommand{\inner}[2]{\left\langle #1,#2 \right\rangle}
\newcommand{\dd}{\,\mathrm{d}}
\newcommand{\law}{\mathrm{Law}}

% ---------- Theorem environments ----------
\newtheorem{theorem}{Theorem}[section]
\newtheorem{proposition}[theorem]{Proposition}
\newtheorem{lemma}[theorem]{Lemma}
\newtheorem{corollary}[theorem]{Corollary}
\theoremstyle{definition}
\newtheorem{definition}[theorem]{Definition}
\newtheorem{assumption}[theorem]{Assumption}
\newtheorem{example}[theorem]{Example}
\theoremstyle{remark}
\newtheorem{remark}[theorem]{Remark}

% ---------- Pastel boxes ----------
\definecolor{pastelblue}{RGB}{233,243,255}
\definecolor{pastelgreen}{RGB}{236,250,241}
\definecolor{pastelpink}{RGB}{255,238,242}
\definecolor{pastelyellow}{RGB}{255,250,230}
\definecolor{pastelgray}{RGB}{245,246,248}
\definecolor{pastelpurple}{RGB}{243,237,255}

\newtcolorbox{keybox}[2][]{
  enhanced,breakable,colback=pastelblue,colframe=blue!15,boxrule=0.6pt,arc=2mm,
  left=2mm,right=2mm,top=1.5mm,bottom=1.5mm,title=\textbf{#2},fonttitle=\normalsize,#1
}
\newtcolorbox{resultbox}[2][]{
  enhanced,breakable,colback=pastelgreen,colframe=green!20,boxrule=0.6pt,arc=2mm,
  left=2mm,right=2mm,top=1.5mm,bottom=1.5mm,title=\textbf{#2},fonttitle=\normalsize,#1
}
\newtcolorbox{warnbox}[2][]{
  enhanced,breakable,colback=pastelyellow,colframe=yellow!35!black,boxrule=0.6pt,arc=2mm,
  left=2mm,right=2mm,top=1.5mm,bottom=1.5mm,title=\textbf{#2},fonttitle=\normalsize,#1
}
\newtcolorbox{insightbox}[2][]{
  enhanced,breakable,colback=pastelpink,colframe=red!15,boxrule=0.6pt,arc=2mm,
  left=2mm,right=2mm,top=1.5mm,bottom=1.5mm,title=\textbf{#2},fonttitle=\normalsize,#1
}
\newtcolorbox{summarybox}[2][]{
  enhanced,breakable,colback=pastelpurple,colframe=purple!15,boxrule=0.6pt,arc=2mm,
  left=2mm,right=2mm,top=1.5mm,bottom=1.5mm,title=\textbf{#2},fonttitle=\normalsize,#1
}

% ---------- Title ----------
\title{\textbf{Flow Matching and Diffusion Models}\\[6pt]
\large From transport PDEs to conditional velocity fields: a rigorous and unified course}
\author{}
\date{}

\begin{document}
\maketitle

\begin{abstract}
This course presents a self-contained, rigorously organized treatment of Flow Matching~(FM) and its relationship to diffusion models.
The logical chain is: (1)~stochastic dynamics and the Fokker--Planck equation; (2)~deterministic dynamics and the continuity equation with its mass conservation criterion; (3)~the SDE/ODE bridge via the probability-flow ODE, reverse-time SDE, and marginal-preserving families; (4)~the FM framework, which learns a velocity field whose ODE flow transports a source distribution to a target, using conditional paths and the marginalization trick for tractable training; (5)~concrete constructions of probability paths and velocity fields, including the probability-path-first (Eulerian) and flow-first (Lagrangian) perspectives; (6)~the full unification showing that Gaussian FM is equivalent to diffusion, that all four prediction parameterizations are interchangeable, and that all affine flows are equivalent up to time reparameterization.
\end{abstract}

\tableofcontents
\newpage

% =====================================================================================
%  PART I: TRANSPORT PDEs AND DYNAMICS
% =====================================================================================

\part{Transport PDEs and Dynamics}

% =====================================================================================
\section{Core objects and notation}\label{sec:notation}

\paragraph{Setting.}
We work in $\R^d$.
All densities are with respect to the Lebesgue measure.
We write $p_t$ for the density of a random variable $X_t$ at time $t$.

\begin{definition}[Score function]
For a smooth, strictly positive density $p$, the \emph{score} is
\[
s(x) \coloneqq \grad_x \log p(x) = \frac{\grad_x p(x)}{p(x)}.
\]
For a time-indexed family $(p_t)_{t\in[0,T]}$, write $s(x,t) = \grad_x \log p_t(x)$.
\end{definition}

\begin{definition}[Probability path]
A \emph{probability path} is a family of densities $(p_t)_{t\in[0,1]}$ with prescribed boundary conditions, typically $p_0 = p_{\mathrm{src}}$ (source) and $p_1 = p_{\mathrm{tgt}}$ (target).
\end{definition}

\begin{keybox}{Central viewpoint}
Both diffusion models and flow matching specify (explicitly or implicitly) a probability path $(p_t)_{t\in[0,1]}$ from a simple source $p_0$ to a data target $p_1$, and then construct dynamics --- deterministic or stochastic --- whose time marginals match this path.
The entire FM--diffusion bridge rests on this principle.
\end{keybox}

% =====================================================================================
\section{Stochastic dynamics and the Fokker--Planck equation}\label{sec:sde-fp}

We begin with the most general setting: stochastic differential equations and the PDE governing the evolution of their marginal densities.

\subsection{Forward SDE}

\begin{definition}[Forward SDE]\label{def:forward-sde}
Let $f:\R^d\times[0,T]\to\R^d$ be a drift and $g:[0,T]\to\R_{\ge 0}$ a scalar diffusion coefficient.
A forward SDE is
\begin{equation}\label{eq:forward-sde}
\dd X_t = f(X_t,t)\dd t + g(t)\dd W_t, \qquad X_0 \sim p_0,
\end{equation}
where $W_t$ is a standard $d$-dimensional Brownian motion.
Under standard regularity (e.g., $f$ Lipschitz in $x$, at most linear growth), \eqref{eq:forward-sde} admits a unique strong solution with continuous sample paths.
\end{definition}

\subsection{The Fokker--Planck equation}\label{subsec:fp}

The PDE governing the evolution of $p_t = \law(X_t)$ is the Fokker--Planck equation.
We state it first in weak form (the most natural formulation from the probabilistic viewpoint), then in strong form.

\begin{definition}[Fokker--Planck equation --- weak form]\label{def:fp-weak}
Let $\varphi \in C_c^\infty(\R^d)$ be a test function.
A family $(p_t)$ satisfies the Fokker--Planck equation for the SDE~\eqref{eq:forward-sde} if, for all $\varphi$,
\begin{equation}\label{eq:fp-weak}
\frac{\dd}{\dd t}\int_{\R^d} \varphi(x)\,p_t(x)\dd x
=
\int_{\R^d} \inner{\grad\varphi(x)}{f(x,t)}\,p_t(x)\dd x
+ \frac{g(t)^2}{2}\int_{\R^d} \Delta\varphi(x)\,p_t(x)\dd x.
\end{equation}
\end{definition}

\begin{proof}[Derivation of the weak form]
Apply It\^o's formula to $\varphi(X_t)$ for the SDE~\eqref{eq:forward-sde}:
\[
\dd\varphi(X_t) = \inner{\grad\varphi(X_t)}{f(X_t,t)}\dd t + g(t)\inner{\grad\varphi(X_t)}{\dd W_t} + \frac{g(t)^2}{2}\Delta\varphi(X_t)\dd t.
\]
Take expectations.
The It\^o integral $\int_0^t g(s)\inner{\grad\varphi(X_s)}{\dd W_s}$ is a martingale (under standard integrability), hence has zero expectation.
Differentiating $\E[\varphi(X_t)] = \int \varphi(x)\,p_t(x)\dd x$ yields~\eqref{eq:fp-weak}.
\end{proof}

\begin{definition}[Fokker--Planck equation --- strong form]\label{def:fp-strong}
If $p_t$ is smooth and decays suitably at infinity, integration by parts in~\eqref{eq:fp-weak} yields the strong form:
\begin{equation}\label{eq:fp-strong}
\partial_t p_t = -\diver(p_t f) + \frac{g(t)^2}{2}\,\Delta p_t.
\end{equation}
\end{definition}

\subsection{A fundamental score identity}\label{subsec:score-id}

The following identity is used repeatedly throughout the course to trade Laplacians for score-weighted divergences.

\begin{lemma}[Score--Laplacian identity]\label{lem:score-lap}
Let $p$ be smooth and strictly positive, with score $s = \grad\log p$.
Then
\begin{equation}\label{eq:score-lap}
\diver(p\,s) = \diver(p\,\grad\log p) = \diver(\grad p) = \Delta p.
\end{equation}
\end{lemma}

\begin{proof}
$p\,\grad\log p = p \cdot \frac{\grad p}{p} = \grad p$.
Hence $\diver(p\,s) = \diver(\grad p) = \Delta p$.
\end{proof}

\begin{insightbox}{Why scores enter all marginal equivalences}
Identity~\eqref{eq:score-lap} implies that the diffusion term $\frac{g^2}{2}\Delta p$ in Fokker--Planck can be rewritten as $\frac{g^2}{2}\diver(p\,s)$.
This is exactly what allows trading stochasticity (diffusion) for a deterministic correction involving the score, producing ODEs and families of SDEs with identical marginals.
\end{insightbox}

% =====================================================================================
\section{Deterministic dynamics and the continuity equation}\label{sec:ode-cont}

We now turn to the deterministic counterpart: ODE flows and the continuity equation.

\subsection{Flow maps and velocity fields}

\begin{definition}[Flow map and pushforward]\label{def:flow-map}
A \emph{flow map} is a family of measurable maps $\Psi_{s\to t}:\R^d\to\R^d$ indexed by $s,t\in[0,1]$, with $\Psi_{s\to s} = \mathrm{Id}$.
Given $X_s \sim p_s$, define $X_t \coloneqq \Psi_{s\to t}(X_s)$.
The law of $X_t$ is the \emph{pushforward} $(\Psi_{s\to t})_\# p_s$, defined by
\[
\int \varphi(x)\,(\Psi_{s\to t})_\# p_s(\dd x) \coloneqq \int \varphi(\Psi_{s\to t}(x_s))\,p_s(\dd x_s), \quad \forall\,\varphi\in C_b(\R^d).
\]
We write $p_t = (\Psi_{s\to t})_\# p_s$ when densities exist.
\end{definition}

\begin{definition}[Velocity field induced by a flow map]\label{def:vel-from-flow}
Assume $\Psi_{0\to t}$ is differentiable in $t$ and invertible for $t\in(0,1)$.
The \emph{velocity field} $v_t:\R^d\to\R^d$ is defined by
\begin{equation}\label{eq:vel-from-flow}
v_t(x) \coloneqq \partial_t \Psi_{0\to t}\!\big(\Psi_{0\to t}^{-1}(x)\big).
\end{equation}
Then trajectories satisfy the ODE
\begin{equation}\label{eq:flow-ode}
\frac{\dd}{\dd t}X_t = v_t(X_t), \qquad X_0 \sim p_0.
\end{equation}
\end{definition}

\begin{warnbox}{Flow map $\longleftrightarrow$ ODE: rigorous correspondence}
If $\Psi_{0\to t}$ is invertible and differentiable, then $v_t$ defined by~\eqref{eq:vel-from-flow} implies $X_t = \Psi_{0\to t}(X_0)$ solves $\dot{X}_t = v_t(X_t)$.
Conversely, if $v_t$ is (globally) Lipschitz in $x$ (uniformly in $t$), then the ODE defines a unique flow map $\Psi_{0\to t}$ by $\Psi_{0\to t}(x_0) = X_t^{x_0}$.
\end{warnbox}

\subsection{The continuity equation}

\begin{definition}[Continuity equation --- weak form]\label{def:cont-weak}
A family of densities $(p_t)$ satisfies the continuity equation with velocity field $v_t$ if, for all $\varphi\in C_c^\infty(\R^d)$,
\begin{equation}\label{eq:cont-weak}
\frac{\dd}{\dd t}\int_{\R^d} \varphi(x)\,p_t(x)\dd x = \int_{\R^d} \inner{\grad\varphi(x)}{v_t(x)}\,p_t(x)\dd x.
\end{equation}
\end{definition}

\begin{proof}[Derivation of the weak form]
If $\dot{X}_t = v_t(X_t)$ and $\varphi\in C_c^\infty$, then by the chain rule:
\[
\frac{\dd}{\dd t}\varphi(X_t) = \inner{\grad\varphi(X_t)}{v_t(X_t)}.
\]
Taking expectations yields~\eqref{eq:cont-weak}.
\end{proof}

\begin{definition}[Continuity equation --- strong form]\label{def:cont-strong}
If $p_t$ is smooth, integration by parts gives the strong form:
\begin{equation}\label{eq:cont-strong}
\partial_t p_t + \diver(p_t\,v_t) = 0.
\end{equation}
\end{definition}

\begin{remark}[Continuity = Fokker--Planck with $g\equiv 0$]
Setting $g(t)\equiv 0$ in the Fokker--Planck equation~\eqref{eq:fp-strong} yields exactly the continuity equation~\eqref{eq:cont-strong} with $v_t = f(\cdot,t)$.
The continuity equation is the deterministic analogue of Fokker--Planck.
\end{remark}

\subsection{Mass conservation criterion}\label{subsec:mass-conservation}

Given a prescribed density path $(p_t)$ and a velocity field $(v_t)$, we want to know: \emph{does the ODE flow induced by $v_t$ produce exactly the densities $p_t$?}

\begin{theorem}[Mass conservation criterion]\label{thm:mass-conservation}
Let $(p_t)_{t\in[0,1]}$ be a prescribed smooth density path and $v_t$ a velocity field defining the ODE $\dot{X}_t = v_t(X_t)$ with flow map $\Psi_{0\to t}$.
Let $p_t^{\mathrm{fwd}} \coloneqq (\Psi_{0\to t})_\# p_0$ be the flow-induced density.
Then, under standard regularity,
\[
p_t^{\mathrm{fwd}} = p_t \quad \text{for all } t\in[0,1]
\]
if and only if the pair $(p_t, v_t)$ satisfies the continuity equation:
\[
\partial_t p_t(x) + \grad_x\cdot\big(p_t(x)\,v_t(x)\big) = 0, \qquad \forall\,t\in[0,1],\ x\in\R^d.
\]
\end{theorem}

\begin{proof}[Proof sketch]
$(\Rightarrow)$: If $p_t^{\mathrm{fwd}} = p_t$, the weak continuity equation holds because, for any test function $\varphi$,
\[
\frac{\dd}{\dd t}\int \varphi(x)\,p_t(x)\dd x = \frac{\dd}{\dd t}\E[\varphi(X_t)] = \E[\inner{\grad\varphi(X_t)}{v_t(X_t)}] = \int \inner{\grad\varphi(x)}{v_t(x)}\,p_t(x)\dd x.
\]
$(\Leftarrow)$: If $(p_t,v_t)$ satisfies the continuity equation, then $p_t$ and $p_t^{\mathrm{fwd}}$ both solve the same linear PDE (the continuity equation) with the same initial condition $p_0$.
Under uniqueness of weak solutions (guaranteed by Lipschitz regularity of $v_t$), they must coincide.
\end{proof}

\begin{keybox}{Why the continuity equation is the underlying rule}
The continuity equation is the single PDE criterion ensuring that a velocity field $v_t$ generates the correct density path $p_t$.
It is the \emph{underlying rule} connecting:
\begin{itemize}[nosep]
\item The probability-flow ODE in diffusion models (Section~\ref{sec:sde-ode-bridge}).
\item The flow matching ODE (Section~\ref{sec:fm-framework}).
\item The conditional-to-marginal construction (Proposition~\ref{prop:marg-vf-generates}).
\end{itemize}
\end{keybox}

\begin{remark}[Non-uniqueness of $v_t$ for a given $p_t$]
The continuity equation $\partial_t p_t + \diver(p_t v_t) = 0$ is one scalar equation for $d$ components of $v_t$: it is underdetermined.
If $v_t$ is a solution and $\tilde{v}_t$ satisfies $\diver(p_t \tilde{v}_t)=0$ (divergence-free with respect to $p_t$), then $v_t + \tilde{v}_t$ also solves the same equation.
Flow matching resolves this non-uniqueness by \emph{constructing} a particular $v_t$ via conditional fields and regression, rather than solving the PDE directly.
\end{remark}

% =====================================================================================
\section{The SDE/ODE bridge}\label{sec:sde-ode-bridge}

We now establish the central PDE-level result: given a stochastic process with known marginals, one can construct deterministic and stochastic dynamics sharing the same marginals, using only the score function.

\subsection{Probability-flow ODE}\label{subsec:pf-ode}

\begin{theorem}[Probability-flow ODE]\label{thm:pf-ode}
Let $(p_t)$ be the marginal densities of the forward SDE~\eqref{eq:forward-sde}, assumed smooth and strictly positive for $t\in(0,T)$.
Define the ODE
\begin{equation}\label{eq:pf-ode}
\frac{\dd}{\dd t}\tilde{X}_t = f(\tilde{X}_t,t) - \frac{g(t)^2}{2}\,\grad_x\log p_t(\tilde{X}_t).
\end{equation}
Then $\tilde{X}_t$ has the same time-marginal densities $p_t$ as the SDE, under standard PDE uniqueness.
\end{theorem}

\begin{proof}
The SDE density satisfies Fokker--Planck:
\[
\partial_t p_t = -\diver(p_t f) + \frac{g(t)^2}{2}\,\Delta p_t.
\]
The ODE~\eqref{eq:pf-ode} has drift $\tilde{f} = f - \frac{g^2}{2}s$ (where $s = \grad\log p_t$), so its density satisfies the continuity equation:
\[
\partial_t \tilde{p}_t = -\diver\!\Big(\tilde{p}_t\big(f - \tfrac{g^2}{2}s\big)\Big) = -\diver(\tilde{p}_t f) + \frac{g^2}{2}\diver(\tilde{p}_t s).
\]
By the score--Laplacian identity~\eqref{eq:score-lap}, $\diver(p_t s) = \Delta p_t$.
Provided $\tilde{p}_t = p_t$, the right-hand side becomes $-\diver(p_t f) + \frac{g^2}{2}\Delta p_t$, which equals $\partial_t p_t$ by Fokker--Planck.
Hence $p_t$ satisfies the continuity equation for $\tilde{f}$, and by uniqueness, $\tilde{p}_t = p_t$.
\end{proof}

\begin{insightbox}{The probability-flow ODE in words}
The PF-ODE removes all stochasticity from the forward SDE while preserving marginal distributions.
The ``missing'' diffusion is exactly compensated by a deterministic score correction $-\frac{g^2}{2}\grad\log p_t$.
This is possible because $\diver(p\,\grad\log p) = \Delta p$: the score encodes precisely the density-level effect of diffusion.
\end{insightbox}

\subsection{Reverse-time SDE}\label{subsec:reverse-sde}

\begin{theorem}[Reverse-time SDE]\label{thm:reverse-sde}
Let $(p_t)$ be the marginal densities of the forward SDE~\eqref{eq:forward-sde} with $X_0 \sim p_{\mathrm{data}}$.
Define the reversed-time density $\tilde{p}_\tau(x) \coloneqq p_{T-\tau}(x)$ for $\tau\in[0,T]$.
The reverse-time SDE
\begin{equation}\label{eq:reverse-sde}
\dd \bar{X}_\tau = \Big[-f(\bar{X}_\tau, T-\tau) + g(T-\tau)^2\,\grad\log p_{T-\tau}(\bar{X}_\tau)\Big]\dd\tau + g(T-\tau)\dd\bar{W}_\tau,
\end{equation}
with $\bar{X}_0 \sim p_T$, has marginal densities $\tilde{p}_\tau = p_{T-\tau}$.
\end{theorem}

\begin{proof}
Since $\tilde{p}_\tau = p_{T-\tau}$, we have $\partial_\tau \tilde{p}_\tau = -\partial_t p_t\big|_{t=T-\tau}$.
Using the forward Fokker--Planck:
\[
\partial_\tau \tilde{p}_\tau = \diver(\tilde{p}_\tau\,f(\cdot,T-\tau)) - \frac{g(T-\tau)^2}{2}\Delta\tilde{p}_\tau.
\]
On the other hand, if~\eqref{eq:reverse-sde} holds with drift $\tilde{h}$, then $\tilde{p}_\tau$ must satisfy
\[
\partial_\tau \tilde{p}_\tau = -\diver(\tilde{p}_\tau\,\tilde{h}) + \frac{g(T-\tau)^2}{2}\Delta\tilde{p}_\tau.
\]
Equating and using $\Delta\tilde{p}_\tau = \diver(\tilde{p}_\tau\grad\log\tilde{p}_\tau)$ yields
\[
\tilde{h}(x,\tau) = -f(x,T-\tau) + g(T-\tau)^2\grad\log p_{T-\tau}(x),
\]
which is the drift in~\eqref{eq:reverse-sde}.
\end{proof}

\begin{remark}
In the original time variable $t$ (decreasing from $T$ to $0$), the reverse SDE reads:
\[
\dd X_t = \big[f(X_t,t) - g(t)^2\,\grad\log p_t(X_t)\big]\dd t + g(t)\dd\bar{W}_t,
\]
where the correction $-g^2 s$ appears with coefficient $1$ (not $\frac{1}{2}$ as in the PF-ODE) because the diffusion term contributes an additional $+\frac{g^2}{2}s$ to the Fokker--Planck balance.
\end{remark}

\subsection{A continuum of marginal-preserving SDEs}\label{subsec:marginal-family}

\begin{theorem}[Marginal-preserving SDE family]\label{thm:marginal-family}
Assume $(p_t)$ is smooth and strictly positive and satisfies the continuity equation with velocity field $v_t$:
\[
\partial_t p_t + \diver(p_t\,v_t) = 0.
\]
Let $\gamma(t)\ge 0$ be any scalar schedule.
Define the SDE
\begin{equation}\label{eq:marginal-family}
\dd X_t = b(X_t,t)\dd t + \gamma(t)\dd B_t, \qquad X_0 \sim p_0,
\end{equation}
with drift
\begin{equation}\label{eq:marginal-family-drift}
b(x,t) \coloneqq v_t(x) + \frac{\gamma(t)^2}{2}\,\grad\log p_t(x).
\end{equation}
Then $(p_t)$ also satisfies the Fokker--Planck equation for~\eqref{eq:marginal-family}.
In particular, $\law(X_t) = p_t$ for all $t$.
\end{theorem}

\begin{proof}
The Fokker--Planck right-hand side for~\eqref{eq:marginal-family} is:
\[
-\diver(p_t b) + \frac{\gamma^2}{2}\Delta p_t
= -\diver(p_t v_t) - \frac{\gamma^2}{2}\diver(p_t\,s) + \frac{\gamma^2}{2}\Delta p_t.
\]
By $\diver(p_t s) = \Delta p_t$, the last two terms cancel, leaving $-\diver(p_t v_t) = \partial_t p_t$.
\end{proof}

\begin{summarybox}{Three key dynamics from one density path}
Given a density path $(p_t)$ with score $s(\cdot,t) = \grad\log p_t$:
\begin{enumerate}[nosep]
\item \textbf{ODE} ($\gamma\equiv 0$): $\dd X_t = v_t(X_t)\dd t$ \hfill (deterministic, PF-ODE)
\item \textbf{Reverse SDE} ($\gamma = g$): the canonical Anderson reversal.
\item \textbf{General SDE} (arbitrary $\gamma$): interpolates between deterministic and stochastic.
\end{enumerate}
All three share the same marginal densities $p_t$, by Theorem~\ref{thm:marginal-family}.
\end{summarybox}

\newpage
% =====================================================================================
%  PART II: FLOW MATCHING FRAMEWORK
% =====================================================================================

\part{Flow Matching Framework}

% =====================================================================================
\section{From diffusion insights to the general framework}\label{sec:fm-framework}

\subsection{Lesson from score-based diffusion}\label{subsec:lesson-diffusion}

Before defining flow matching in full generality, we extract the key structural insights from diffusion models (Chapter~\ref{sec:sde-ode-bridge}).
The construction of a diffusion-based generative model follows three steps:

\paragraph{Step 1: A conditional path and its marginal densities.}
A diffusion model defines a family of densities $(p_t)_{t\in[0,1]}$ via a forward conditional distribution
\begin{equation}\label{eq:diff-cond-path}
p_t(x_t\mid x_0) = \cN(x_t;\,\alpha_t x_0,\,\sigma_t^2 \Id), \qquad x_0\sim p_{\mathrm{data}},
\end{equation}
which induces the marginal density $p_t(x_t) = \int p_t(x_t\mid x_0)\,p_{\mathrm{data}}(x_0)\dd x_0$.
The increasing variance $\sigma_t^2$ drives $p_t$ toward a Gaussian prior.

\paragraph{Step 2: A velocity field from the Fokker--Planck equation.}
The marginal $p_t$ evolves according to a velocity field $v_t(x) = f(t)x - \frac{1}{2}g^2(t)\grad\log p_t(x)$, derived from the probability-flow ODE (Theorem~\ref{thm:pf-ode}).
The ODE $\dot{x}(t) = v_t(x(t))$ transports samples along $(p_t)$.
The underlying PDE ensuring this alignment is the continuity equation (Theorem~\ref{thm:mass-conservation}).

\paragraph{Step 3: Learning via the conditional strategy.}
Since $\grad\log p_t(x_t)$ is intractable, we exploit the tractable conditional score $\grad\log p_t(x_t\mid x_0)$.
The key identity
\[
\grad\log p_t(x_t) = \E_{x_0\sim p(\cdot\mid x_t)}\!\big[\grad\log p_t(x_t\mid x_0)\big]
\]
shows that marginal scores are conditional expectations of conditional scores.
Training with conditional targets (denoising score matching) differs from the marginal objective only by a $\theta$-independent constant.

\begin{keybox}{The three-step pattern}
\begin{enumerate}[nosep]
\item \textbf{Define} a conditional path $p_t(\cdot\mid z)$ and its marginal $p_t = \int p_t(\cdot\mid z)\pi(z)\dd z$.
\item \textbf{Identify} a velocity field $v_t$ such that $(p_t, v_t)$ satisfies the continuity equation.
\item \textbf{Learn} $v_t$ by regressing against tractable conditional targets $v_t(\cdot\mid z)$.
\end{enumerate}
Flow matching generalizes this pattern to \emph{arbitrary} source and target distributions.
\end{keybox}

\subsection{The flow matching framework}\label{subsec:fm-general}

We now present the FM framework in its full generality, following the same three-step structure.

\paragraph{Setup.}
Let $p_{\mathrm{src}}$ and $p_{\mathrm{tgt}}$ be two distributions on $\R^d$ (source and target), both easy to sample from.
We set $p_0 = p_{\mathrm{src}}$ and $p_1 = p_{\mathrm{tgt}}$.
The generation task (source = Gaussian prior, target = data) is a special case.

\subsubsection{Step 1: Conditional path and its marginal densities}

FM implicitly defines a continuous family $(p_t)_{t\in[0,1]}$ interpolating between the endpoints.
Each marginal $p_t$ is expressed via a latent variable $z\sim\pi(z)$ and a conditional distribution $p_t(\cdot\mid z)$:
\begin{equation}\label{eq:fm-marginal}
p_t(x_t) = \int p_t(x_t\mid z)\,\pi(z)\dd z,
\end{equation}
with $(\pi, \{p_t(\cdot\mid z)\})$ chosen to satisfy the boundary conditions $p_0 = p_{\mathrm{src}}$, $p_1 = p_{\mathrm{tgt}}$.

Common choices for the latent variable $z$:
\begin{itemize}[nosep]
\item \textbf{Two-sided conditioning}: $z = (x_0, x_1)$ with $\pi = p_{\mathrm{src}}(x_0)\,p_{\mathrm{tgt}}(x_1)$.
\item \textbf{One-sided conditioning}: $z = x_1 \sim p_{\mathrm{tgt}}$ (or $z = x_0\sim p_{\mathrm{src}}$), recovering diffusion-like setups when $p_{\mathrm{src}}$ is Gaussian.
\end{itemize}

\subsubsection{Step 2: Velocity field}

The goal is to find a velocity field $v_t(x)$ such that the ODE
\begin{equation}\label{eq:fm-ode}
\frac{\dd x(t)}{\dd t} = v_t(x(t)), \qquad x(0)\sim p_{\mathrm{src}},
\end{equation}
produces marginal distributions matching $p_t$ at each time $t$.
By Theorem~\ref{thm:mass-conservation}, this requires $(p_t, v_t)$ to satisfy the continuity equation:
\begin{equation}\label{eq:fm-cont}
\partial_t p_t(x) + \grad_x\cdot\big(p_t(x)\,v_t(x)\big) = 0.
\end{equation}
As noted in Section~\ref{subsec:mass-conservation}, $v_t$ is not unique for a given $p_t$.
FM seeks one particular solution by marginalization of conditional velocity fields.

\subsubsection{Step 3: Learning via the conditional strategy}

\paragraph{Conditional velocity field.}
For each $z$, we define (or derive) a conditional velocity field $v_t(\cdot\mid z)$ such that the conditional ODE
\[
\dot{X}_t = v_t(X_t\mid z), \qquad X_0\sim p_0(\cdot\mid z),
\]
generates $p_t(\cdot\mid z)$.

\paragraph{Marginal velocity by marginalization.}
The marginal velocity field is defined as
\begin{equation}\label{eq:marg-vf}
v_t(x) \coloneqq \int v_t(x\mid z)\,\frac{p_t(x\mid z)\,\pi(z)}{p_t(x)}\dd z = \E_{z\sim p(\cdot\mid x_t=x)}\!\big[v_t(x\mid z)\big].
\end{equation}

\paragraph{FM and CFM losses.}
The FM loss is
\[
\mathcal{L}_{\mathrm{FM}}(\theta) = \E_{t,\,x_t\sim p_t}\!\big[\norm{v_\theta(x_t,t) - v_t(x_t)}^2\big],
\]
which is intractable because $v_t$ involves an expectation over the posterior $p(z\mid x_t)$.
The \emph{conditional} FM loss is
\begin{equation}\label{eq:cfm-loss}
\mathcal{L}_{\mathrm{CFM}}(\theta) = \E_{t,\,z\sim\pi,\,x_t\sim p_t(\cdot\mid z)}\!\big[\norm{v_\theta(x_t,t) - v_t(x_t\mid z)}^2\big].
\end{equation}
We prove below that these two losses have the same gradient and the same minimizer.

\subsection{The marginalization trick}\label{subsec:marg-trick}

\begin{proposition}[Marginal velocity field generates the marginal path]\label{prop:marg-vf-generates}
If each conditional velocity field $v_t(\cdot\mid z)$ generates $p_t(\cdot\mid z)$ (i.e., satisfies $\partial_t p_t(\cdot\mid z) + \diver(p_t(\cdot\mid z)\,v_t(\cdot\mid z)) = 0$), then the marginal velocity field $v_t$ defined by~\eqref{eq:marg-vf} generates the marginal path $p_t$.
\end{proposition}

\begin{proof}
Fix $\varphi\in C_c^\infty(\R^d)$.
For each $z$, the conditional continuity equation in weak form gives:
\[
\frac{\dd}{\dd t}\int \varphi(x)\,p_t(x\mid z)\dd x = \int \inner{\grad\varphi(x)}{v_t(x\mid z)}\,p_t(x\mid z)\dd x.
\]
Multiply by $\pi(z)$ and integrate over $z$:
\[
\frac{\dd}{\dd t}\int \varphi(x)\underbrace{\Big(\int p_t(x\mid z)\pi(z)\dd z\Big)}_{p_t(x)}\dd x
= \int \inner{\grad\varphi(x)}{\underbrace{\int v_t(x\mid z)\,p_t(x\mid z)\,\pi(z)\dd z}_{v_t(x)\,p_t(x)}}\dd x.
\]
This is exactly the weak continuity equation for $(p_t, v_t)$.
By Theorem~\ref{thm:mass-conservation}, the ODE flow of $v_t$ generates $p_t$.
\end{proof}

\begin{keybox}{The marginalization trick in one line}
\[
v_t(x) = \E\!\big[v_t(x\mid Z)\,\big|\,X_t = x\big].
\]
This identity converts a family of tractable conditional velocity fields into a single marginal velocity field that generates the correct marginal density path.
\end{keybox}

\subsection{CFM and FM have the same gradient}\label{subsec:cfm-fm-equiv}

\begin{lemma}[Variance decomposition for regression]\label{lem:var-decomp}
Let $A$ be square-integrable and $X$ any random variable.
For any measurable predictor $f(X)$,
\[
\E\!\big[\norm{f(X)-A}^2\big] = \E\!\big[\norm{f(X)-\E[A\mid X]}^2\big] + \E\!\big[\norm{A-\E[A\mid X]}^2\big].
\]
\end{lemma}

\begin{proof}
Expand $\norm{f-A}^2 = \norm{(f-\E[A\mid X]) + (\E[A\mid X]-A)}^2$ and note that the cross term has zero expectation because $\E[A - \E[A\mid X]\mid X] = 0$.
\end{proof}

\begin{theorem}[Equivalence of CFM and FM]\label{thm:cfm-fm}
Let $v_t(x) = \E[v_t(x\mid Z)\mid X_t = x]$ be the marginal velocity.
Then
\[
\mathcal{L}_{\mathrm{CFM}}(\theta) = \mathcal{L}_{\mathrm{FM}}(\theta) + C,
\]
where $C$ does not depend on $\theta$.
In particular, $\nabla_\theta\mathcal{L}_{\mathrm{CFM}} = \nabla_\theta\mathcal{L}_{\mathrm{FM}}$, and the unique minimizer satisfies
\[
v^*_\theta(x_t,t) = v_t(x_t) = \E\!\big[v_t(x_t\mid Z)\,\big|\,X_t = x_t\big].
\]
\end{theorem}

\begin{proof}
Apply Lemma~\ref{lem:var-decomp} with $A = v_t(X_t\mid Z)$, $X = X_t$, and $f = v_\theta(\cdot,t)$.
Then $\E[A\mid X_t] = v_t(X_t)$ by definition~\eqref{eq:marg-vf}, and the residual term $\E[\norm{A - \E[A\mid X_t]}^2]$ is $\theta$-independent.
\end{proof}

\begin{resultbox}{Practical meaning}
Training with tractable conditional targets $v_t(\cdot\mid z)$ (CFM) is not an approximation.
It is an exactly justified surrogate: same gradients, same minimizer as training with the intractable marginal target $v_t$.
This is the same principle as denoising score matching in diffusion models.
\end{resultbox}

\newpage
% =====================================================================================
%  PART III: CONSTRUCTING PATHS AND VELOCITIES
% =====================================================================================

\part{Constructing Probability Paths and Velocity Fields}

% =====================================================================================
\section{Warm-up: Gaussian-to-Gaussian bridge}\label{sec:gauss-bridge}

When both $p_{\mathrm{src}}$ and $p_{\mathrm{tgt}}$ are Gaussian, the marginal density and velocity can be computed analytically.
This serves as a template for the general conditional construction.

Consider the interpolated Gaussian path:
\begin{equation}\label{eq:gauss-path}
p_t(x) = \cN\!\big(x;\,\mu(t),\,\sigma(t)^2 \Id\big),
\end{equation}
with endpoints $p_0 = \cN(\mu(0), \sigma(0)^2\Id)$ and $p_1 = \cN(\mu(1), \sigma(1)^2\Id)$.

A natural affine flow map realizing this path is:
\begin{equation}\label{eq:gauss-flow}
\Psi_{0\to t}(x) = \mu(t) + \frac{\sigma(t)}{\sigma(0)}\big(x - \mu(0)\big).
\end{equation}
If $x\sim p_0$, then $\Psi_{0\to t}(x) \sim \cN(\mu(t), \sigma(t)^2\Id) = p_t$.

\begin{proposition}[Velocity field for a Gaussian density path]\label{prop:gauss-velocity}
The velocity field generating the flow~\eqref{eq:gauss-flow} is:
\begin{equation}\label{eq:gauss-velocity}
v_t(x) = \frac{\sigma'(t)}{\sigma(t)}\big(x - \mu(t)\big) + \mu'(t).
\end{equation}
\end{proposition}

\begin{proof}
Differentiate $\Psi_{0\to t}(y)$ in $t$:
$\Psi'_{0\to t}(y) = \mu'(t) + \frac{\sigma'(t)}{\sigma(0)}(y - \mu(0))$.
The inverse is $y = \Psi_{t\to 0}(x) = \mu(0) + \frac{\sigma(0)}{\sigma(t)}(x - \mu(t))$.
Substituting:
$v_t(x) = \Psi'_{0\to t}(\Psi_{t\to 0}(x)) = \mu'(t) + \frac{\sigma'(t)}{\sigma(t)}(x - \mu(t))$.
\end{proof}

\begin{remark}[Two perspectives on uniqueness]
For a \emph{fixed flow map} $\Psi_{0\to t}$, the velocity is uniquely determined by $v_t = \partial_t\Psi_{0\to t}\circ\Psi_{0\to t}^{-1}$.
For a \emph{fixed density path} $p_t$ (without fixing $\Psi_{0\to t}$), the velocity is not unique: the continuity equation admits infinitely many solutions.
This distinction precisely characterizes the difference between the flow-first and probability-path-first perspectives below.
\end{remark}

% =====================================================================================
\section{Probability-path-first construction (Eulerian view)}\label{sec:path-first}

\textbf{Principle:} Choose a conditional probability path $p_t(\cdot\mid z)$ first, then derive the conditional velocity $v_t(\cdot\mid z)$.

\subsection{One-sided conditioning ($z = x_1 \sim p_{\mathrm{tgt}}$)}

Fix $x_1\sim p_{\mathrm{tgt}}$.
Define the conditional Gaussian path:
\begin{equation}\label{eq:one-sided-path}
p_t(x_t\mid x_1) = \cN\!\big(x_t;\,b_t x_1,\,a_t^2 \Id\big),
\end{equation}
with boundary conditions $a_0 = 1,\,b_0 = 0$ and $a_1 = 0,\,b_1 = 1$ (the limit $a_1\to 0$ concentrates $p_1(\cdot\mid x_1)$ at $\delta_{x_1}$).

At the boundaries: $p_0(\cdot\mid x_1) = \cN(0,\Id)$ (independent of $x_1$) and $p_1(\cdot\mid x_1) = \delta(\cdot - x_1)$.
Marginalizing: $p_0 = \cN(0,\Id) = p_{\mathrm{src}}$ and $p_1 = p_{\mathrm{tgt}}$.

Applying Proposition~\ref{prop:gauss-velocity} to the conditional Gaussian path~\eqref{eq:one-sided-path} (with $\mu(t) = b_t x_1$ and $\sigma(t) = a_t$):

\begin{proposition}[Conditional velocity for one-sided path-first construction]\label{prop:one-sided-vf}
For $t\in(0,1)$ with $a_t > 0$:
\begin{equation}\label{eq:one-sided-vf}
v_t(x\mid x_1) = b_t' x_1 + \frac{a_t'}{a_t}(x - b_t x_1).
\end{equation}
\end{proposition}

\paragraph{CFM loss.}
Sampling $x_t = a_t\varepsilon + b_t x_1$ with $\varepsilon\sim\cN(0,\Id)$, the CFM loss becomes:
\begin{equation}\label{eq:one-sided-cfm}
\mathcal{L}_{\mathrm{CFM}} = \E_{t,\,x_1\sim p_{\mathrm{tgt}},\,x_t\sim p_t(\cdot\mid x_1)}\!\bigg[\bigg\|v_\theta(x_t,t) - \Big[b_t' x_1 + \frac{a_t'}{a_t}(x_t - b_t x_1)\Big]\bigg\|^2\bigg].
\end{equation}

\paragraph{On paired samples.}
Writing $x_0 = \varepsilon$ and $x_t = a_t x_0 + b_t x_1$, the conditional velocity simplifies to $v_t(x_t\mid x_1) = a_t' x_0 + b_t' x_1$, which is the time derivative of the interpolation $x_t$.

\subsection{Two-sided conditioning ($z = (x_0, x_1)$)}

Let $z = (x_0, x_1)$ with $x_0\sim p_{\mathrm{src}}$, $x_1\sim p_{\mathrm{tgt}}$ independently.
Define:
\begin{equation}\label{eq:two-sided-path}
p_t(x_t\mid x_0, x_1) = \cN\!\big(x_t;\,a_t x_0 + b_t x_1,\,\sigma^2 \Id\big),
\end{equation}
with $a_0=1,\,b_0=0,\,a_1=0,\,b_1=1$.
In the deterministic limit $\sigma=0$: $p_t(\cdot\mid x_0,x_1) = \delta(\cdot - (a_t x_0 + b_t x_1))$.

\begin{proposition}[Conditional velocity for two-sided construction]
\begin{equation}\label{eq:two-sided-vf}
v_t(x\mid x_0,x_1) = a_t' x_0 + b_t' x_1.
\end{equation}
\end{proposition}

The CFM loss (with $\sigma=0$, so $x_t = a_t x_0 + b_t x_1$):
\[
\mathcal{L}_{\mathrm{CFM}} = \E_{t,\,x_0\sim p_{\mathrm{src}},\,x_1\sim p_{\mathrm{tgt}}}\!\big[\norm{v_\theta(x_t,t) - (a_t' x_0 + b_t' x_1)}^2\big].
\]

\begin{remark}[Same minimizer for one-sided and two-sided]
On paired samples $(x_0, x_1)$ with $x_t = a_t x_0 + b_t x_1$, both the one-sided target $v_t(x_t\mid x_1) = a_t' x_0 + b_t' x_1$ and the two-sided target coincide.
The optimal velocity field is:
\[
v^*(x,t) = \E\!\big[a_t' X_0 + b_t' X_1\,\big|\,X_t = x\big].
\]
\end{remark}

\subsection{Gaussian FM = Diffusion}

When $p_{\mathrm{src}} = \cN(0,\Id)$, the one-sided conditioning on $z = x_1\sim p_{\mathrm{data}}$ with Gaussian conditional paths is called \emph{Gaussian FM}.
Under the relabeling $a_t \mapsto \sigma_t$, $b_t \mapsto \alpha_t$ (matching diffusion notation), this is exactly the diffusion model trained to predict the velocity.

The canonical choice $a_t = 1-t$, $b_t = t$ (equivalently $\sigma_t = 1-t$, $\alpha_t = t$) gives the FM/Rectified Flow schedule with CFM target $v_t(x_t\mid x_1) = x_1 - x_0$.

\begin{table}[h]
\centering
\renewcommand{\arraystretch}{1.3}
\begin{tabular}{@{} l c c c c @{}}
\toprule
& \textbf{VE} & \textbf{VP} & \textbf{FM/RF} & \textbf{Trigonometric} \\
\midrule
$a_t$ (source coeff.) & $a_t$ & $\sqrt{1-b_t^2}$ & $1-t$ & $\cos(\frac{\pi}{2}t)$ \\
$b_t$ (target coeff.) & $1$ & $b_t$ & $t$ & $\sin(\frac{\pi}{2}t)$ \\
$a_0,\,b_0$ & $\cdot,\,1$ & $\cdot,\,1$ & $1,\,0$ & $1,\,0$ \\
$a_1,\,b_1$ & $a_1,\,1$ & $0,\,\cdot$ & $0,\,1$ & $0,\,1$ \\
$p_{\mathrm{src}}$ & $\cN(0,a_1^2\Id)$ & $\cN(0,\Id)$ & $\cN(0,\Id)$ & $\cN(0,\Id)$ \\
\bottomrule
\end{tabular}
\caption{Common interpolant schedules in the FM convention $x_t = a_t x_0 + b_t x_1$, where $x_0\sim p_{\mathrm{src}}$ and $x_1\sim p_{\mathrm{tgt}}$.}
\label{tab:schedules}
\end{table}

% =====================================================================================
\section{Flow-first construction (Lagrangian view)}\label{sec:flow-first}

\textbf{Principle:} Specify a conditional flow map $\Psi_{0\to t}(\cdot\mid z)$ directly, then derive the conditional velocity by time-differentiation along trajectories.

\subsection{Conditional affine flow}

Fix $z\sim\pi$ and define the time-varying conditional affine flow:
\begin{equation}\label{eq:cond-affine-flow}
\Psi_{0\to t}(x_0\mid z) \coloneqq \mu_t(z) + A_t(z)\,x_0,
\end{equation}
where $\mu_t(z)\in\R^d$ and $A_t(z)\in\R^{d\times d}$ is invertible for $t\in(0,1)$.
Boundary: $A_0(z) = \Id$, $\mu_0(z) = 0$ (recovers $p_{\mathrm{src}}$ at $t=0$).

\paragraph{Induced conditional path.}
$p_t(\cdot\mid z) = (\Psi_{0\to t}(\cdot\mid z))_\# p_{\mathrm{src}}$.
Sampling: draw $x_0\sim p_{\mathrm{src}}$, set $x_t = \mu_t(z) + A_t(z)\,x_0$.

\paragraph{Derived conditional velocity.}
The conditional velocity is obtained by differentiating the flow map along its trajectory:
\begin{equation}\label{eq:flow-first-vf}
v_t(x\mid z) = \frac{\dd}{\dd t}\Psi_{0\to t}\!\big(\Psi_{t\to 0}(x\mid z)\mid z\big).
\end{equation}
Since $x = \mu_t(z) + A_t(z)\,x_0$, we have $x_0 = A_t(z)^{-1}(x - \mu_t(z))$, giving:
\begin{equation}\label{eq:cond-affine-vf}
v_t(x\mid z) = \mu_t'(z) + A_t'(z)\,A_t(z)^{-1}\big(x - \mu_t(z)\big).
\end{equation}

\subsection{One-sided ($z = x_1$)}

Set $\mu_t(z) = b_t x_1$ and $A_t(z) = a_t \Id$ with $a_0=1$, $b_0=0$, $a_1=0$, $b_1=1$:
\[
x_t = a_t x_0 + b_t x_1, \qquad v_t(x\mid x_1) = b_t' x_1 + \frac{a_t'}{a_t}(x - b_t x_1).
\]
On paired samples: $v_t(x_t\mid x_1) = a_t' x_0 + b_t' x_1$, recovering the same CFM target as the path-first approach.

\subsection{Two-sided ($z = (x_0, x_1)$)}

With $\mu_t(x_0,x_1) = b_t x_1$ and $A_t = a_t \Id$:
\[
x_t = a_t x_0 + b_t x_1, \qquad v_t(x_t\mid x_0,x_1) = a_t' x_0 + b_t' x_1.
\]

\subsection{Recovery of the Gaussian bridge}

If $\mu_t$ is independent of $z$ (denoted $\mu(t)$) and $A_t = \frac{\sigma(t)}{\sigma(0)}\Id$, then
\[
\Psi_{0\to t}(x_0) = \mu(t) + \frac{\sigma(t)}{\sigma(0)}(x_0 - \mu(0) + \mu(0)) = \mu(t) + \frac{\sigma(t)}{\sigma(0)}(x_0 - \mu(0)),
\]
and $v_t(x) = \mu'(t) + \frac{\sigma'(t)}{\sigma(t)}(x - \mu(t))$, recovering the Gaussian bridge velocity of Proposition~\ref{prop:gauss-velocity}.

% =====================================================================================
\section{Path-first vs.\ flow-first: comparison}\label{sec:comparison}

\begin{table}[h]
\centering
\renewcommand{\arraystretch}{1.4}
\small
\begin{tabular}{@{} >{\bfseries}p{2.5cm} p{5.5cm} p{5.5cm} @{}}
\toprule
Axis & \textbf{Path-First (Eulerian)} & \textbf{Flow-First (Lagrangian)} \\
\midrule
Given & Conditional density path $p_t(\cdot\mid z)$ & Conditional flow map $\Psi_{0\to t}(\cdot\mid z)$ \\
\addlinespace
Derive & Velocity $v_t(\cdot\mid z)$ by solving the continuity equation for $p_t(\cdot\mid z)$ & Velocity $v_t(x\mid z) = \partial_t\Psi_{0\to t}(\Psi_{t\to 0}(x\mid z)\mid z)$ \\
\addlinespace
Uniqueness of $v_t$ & Not unique for a given $p_t(\cdot\mid z)$ (one scalar PDE, $d$ unknowns) & Unique given $\Psi_{0\to t}(\cdot\mid z)$ \\
\addlinespace
Realizability & Must verify the constructed $v_t$ satisfies the continuity equation & Holds by construction: $p_t(\cdot\mid z) = (\Psi_{0\to t})_\# p_0(\cdot\mid z)$ \\
\addlinespace
Closed form & Convenient when $p_t(\cdot\mid z)$ is Gaussian / exponential family & Convenient when $\Psi_{0\to t}$ has structure (affine, geodesic) \\
\addlinespace
Preferred use & Diffusion-style constructions; analytic conditional Gaussians & Structural priors on trajectories; boundary control; manifolds \\
\bottomrule
\end{tabular}
\caption{Comparison of the two complementary perspectives for constructing conditional paths and velocities.}
\label{tab:path-vs-flow}
\end{table}

\begin{warnbox}{Takeaway}
Path-first is natural when conditional paths admit closed-form velocities (e.g., Gaussians).
Flow-first is natural when you have strong structural priors on trajectories (e.g., affine maps, geodesics on manifolds).
Both yield equivalent transport under regularity and lead to the same CFM training objective.
\end{warnbox}

\newpage
% =====================================================================================
%  PART IV: GAUSSIAN PATHS, SCORES, AND FULL UNIFICATION
% =====================================================================================

\part{Gaussian Paths, Scores, and Full Unification}

In Parts~II and~III we developed the FM framework in full generality: conditional paths, conditional velocity fields, the marginalization trick, and CFM training.
We now specialize to \emph{affine Gaussian} paths, the setting where FM and diffusion models coincide.
The goal of this part is threefold:
\begin{enumerate}[nosep]
\item Show how forward diffusion SDEs produce the affine Gaussian conditional structure.
\item Derive, in a single pass, all score-based identities and the algebraic equivalences among the four prediction types ($\varepsilon$, $x$, score, $v$).
\item State the FM~=~Diffusion equivalence as a clean theorem.
\end{enumerate}

% =====================================================================================
\section{From forward SDEs to affine Gaussian schedules}\label{sec:linear-sde}

This section answers a natural question: \emph{where do the schedules $(\alpha_t,\sigma_t)$ come from?}
The answer: they are the mean and noise coefficients of the Gaussian conditional law of a linear forward SDE.

\subsection{Linear SDE and its Gaussian conditional law}

Consider the linear SDE:
\begin{equation}\label{eq:linear-sde}
\dd Z_t = f_t Z_t\dd t + g_t\dd B_t, \qquad Z_0\sim p_{\mathrm{data}},\qquad t\in[0,T],
\end{equation}
with scalar $f_t\in\R$ and $g_t\ge 0$.

\begin{definition}[Mean coefficient and noise scale]\label{def:alpha-sigma-sde}
\begin{equation}\label{eq:alpha-sde}
\alpha_t \coloneqq \exp\!\Big(\int_0^t f_s\dd s\Big), \qquad
\sigma_t^2 \coloneqq \alpha_t^2\int_0^t g_s^2\,\alpha_s^{-2}\dd s.
\end{equation}
\end{definition}

\begin{lemma}[Gaussian conditional law]\label{lem:linear-sde-gauss}
The unique strong solution of~\eqref{eq:linear-sde} is $Z_t = \alpha_t Z_0 + \alpha_t\int_0^t \alpha_s^{-1}g_s\dd B_s$.
Conditionally on $Z_0$:
\begin{equation}\label{eq:linear-cond}
Z_t\mid Z_0 \;\stackrel{d}{=}\; \alpha_t Z_0 + \sigma_t\varepsilon, \qquad \varepsilon\sim\cN(0,\Id).
\end{equation}
Moreover, $\sigma_t^2$ satisfies $\dot\sigma_t^2 = 2f_t\sigma_t^2 + g_t^2$, $\sigma_0^2 = 0$.
\end{lemma}

\begin{proof}
Apply It\^o to $\alpha_t^{-1}Z_t$ using $\dot\alpha_t = f_t\alpha_t$:
$\dd(\alpha_t^{-1}Z_t) = \alpha_t^{-1}g_t\dd B_t$.
Integrate and use It\^o isometry to compute the variance.
\end{proof}

\begin{keybox}{The Gaussian conditional structure}
Every linear SDE produces Gaussian conditionals $Z_t\mid Z_0 \sim \cN(\alpha_t Z_0,\,\sigma_t^2\Id)$.
This is the same affine Gaussian structure that FM uses for conditional paths (Section~\ref{sec:path-first}).
The forward diffusion SDE is thus one \emph{origin story} for the schedules $(\alpha_t,\sigma_t)$; FM can also choose them freely.
\end{keybox}

\subsection{Inverse identification: $(\alpha_t,\sigma_t) \to (f_t,g_t)$}

\begin{proposition}\label{prop:inverse-id}
If $\alpha_t > 0$ and $\sigma_t > 0$ are differentiable:
\begin{equation}\label{eq:f-from-alpha}
f_t = \frac{\dot\alpha_t}{\alpha_t}, \qquad g_t^2 = \dot\sigma_t^2 - 2f_t\sigma_t^2 = 2\alpha_t\sigma_t\,\frac{\dd}{\dd t}\!\Big(\frac{\sigma_t}{\alpha_t}\Big).
\end{equation}
\end{proposition}

This bijection $(\alpha_t,\sigma_t)\leftrightarrow(f_t,g_t)$ allows seamless translation between FM schedules and SDE coefficients.

\subsection{VP and VE as special cases}

\begin{example}[Variance Preserving (VP)]
Forward SDE: $\dd Y_t = -\frac{1}{2}\beta(t)Y_t\dd t + \sqrt{\beta(t)}\dd W_t$, $Y_0\sim p_{\mathrm{data}}$.
Let $\Gamma(t) = \int_0^t\beta(u)\dd u$. Then $f_t = -\frac{1}{2}\beta(t)$, $g_t = \sqrt{\beta(t)}$, giving:
\[
\alpha_t = e^{-\Gamma(t)/2}, \qquad \sigma_t^2 = 1 - e^{-\Gamma(t)}.
\]
Note $\alpha_t^2 + \sigma_t^2 = 1$: this is the ``variance preserving'' property.
\end{example}

\begin{example}[Variance Exploding (VE)]
Forward SDE: $\dd Y_t = g(t)\dd W_t$, $Y_0\sim p_{\mathrm{data}}$.
Here $f_t = 0$, $\alpha_t \equiv 1$, $\sigma_t^2 = \int_0^t g(u)^2\dd u$.
\end{example}

% =====================================================================================
\section{Score identities, four parameterizations, and the FM = Diffusion equivalence}\label{sec:four-param}

We now work under the affine Gaussian path and derive, in a single unified pass, all score-based identities and prediction equivalences.
This section combines the score/posterior analysis with the four-parameterization framework, avoiding any duplication.

\subsection{The affine Gaussian path}\label{subsec:affine-setup}

\begin{assumption}[Affine Gaussian path --- noise-to-data orientation]\label{ass:affine-gauss}
Let $X_0\sim\cN(0,\Id)$ (noise) and $Z\sim p_{\mathrm{data}}$ (data) be independent.
Let $\alpha_t, \sigma_t$ be $C^1$ schedules on $[0,1]$ with
\[
\alpha_0 = 0,\;\sigma_0 = 1, \qquad \alpha_1 = 1,\;\sigma_1 = 0, \qquad \sigma_t > 0\ \text{for }t\in[0,1).
\]
Define
\begin{equation}\label{eq:affine-path}
\bar{X}_t \coloneqq \sigma_t X_0 + \alpha_t Z,
\end{equation}
with density $p_t$ and score $s(x,t) = \grad_x\log p_t(x)$ for $t\in(0,1)$.
\end{assumption}

Conditionally on $Z = z$, $\bar{X}_t \sim \cN(\alpha_t z,\,\sigma_t^2\Id)$.
This is the same Gaussian conditional structure as~\eqref{eq:linear-cond}, viewed in noise-to-data time.
Both FM and diffusion operate on this shared structure.

\subsection{Score marginalization for Gaussian mixtures}

\begin{lemma}[Score marginalization]\label{lem:score-marg}
Under Assumption~\ref{ass:affine-gauss}, for $t\in(0,1)$:
\begin{equation}\label{eq:score-marg}
\grad_x\log p_t(x) = \E\!\big[\grad_x\log p_t(x\mid Z)\,\big|\,\bar{X}_t = x\big].
\end{equation}
\end{lemma}

\begin{proof}
From $p_t(x) = \int p_t(x\mid z)\,p_{\mathrm{data}}(z)\dd z$, differentiate under the integral:
\[
\grad p_t(x) = \int p_t(x\mid z)\,\grad\log p_t(x\mid z)\,p_{\mathrm{data}}(z)\dd z.
\]
Divide by $p_t(x)$ and recognize the posterior weighting $p(z\mid x) = p_t(x\mid z)\,p_{\mathrm{data}}(z)/p_t(x)$.
\end{proof}

This is the score-level analogue of the marginalization trick (Proposition~\ref{prop:marg-vf-generates}): the marginal score is a posterior expectation of conditional scores.

\subsection{The four prediction types and their algebraic equivalence}

Throughout, we write $x_t = \alpha_t x_0 + \sigma_t\varepsilon$ with $x_0\sim p_{\mathrm{data}}$, $\varepsilon\sim\cN(0,\Id)$, and $p_t(x_t\mid x_0) = \cN(x_t;\,\alpha_t x_0,\,\sigma_t^2\Id)$.

\paragraph{The four predictions.}
All four aim to approximate a conditional expectation given the noisy observation $x_t$:
\begin{enumerate}[nosep,leftmargin=18pt]
\item \textbf{$\varepsilon$-prediction} (noise): $\varepsilon^*(x_t,t) = \E[\varepsilon\mid x_t]$.
\item \textbf{$x$-prediction} (clean): $x^*(x_t,t) = \E[x_0\mid x_t]$.
\item \textbf{Score prediction}: $s^*(x_t,t) = \grad\log p_t(x_t)$.
\item \textbf{$v$-prediction} (velocity): $v^*(x_t,t) = \E[\dot{x}_t\mid x_t]$, where $\dot{x}_t = \alpha_t' x_0 + \sigma_t'\varepsilon$.
\end{enumerate}

\begin{proposition}[Algebraic equivalence of the four predictions]\label{prop:four-equiv}
The four optimal predictions are related by:
\begin{align}
\varepsilon^*(x_t,t) &= -\sigma_t\,s^*(x_t,t), \label{eq:eps-score}\\
x^*(x_t,t) &= \frac{1}{\alpha_t}\big(x_t + \sigma_t^2\,s^*(x_t,t)\big), \label{eq:x-score}\\
v^*(x_t,t) &= \alpha_t'\,x^*(x_t,t) + \sigma_t'\,\varepsilon^*(x_t,t). \label{eq:v-xeps}
\end{align}
\end{proposition}

\begin{proof}
\emph{Equation~\eqref{eq:eps-score}.}
The conditional score is $\grad\log p_t(x_t\mid x_0) = -\frac{x_t - \alpha_t x_0}{\sigma_t^2} = -\frac{\varepsilon}{\sigma_t}$.
By Lemma~\ref{lem:score-marg}:
\[
s^*(x_t,t) = \E\!\Big[-\frac{\varepsilon}{\sigma_t}\,\Big|\,x_t\Big] = -\frac{1}{\sigma_t}\varepsilon^*(x_t,t).
\]

\emph{Equation~\eqref{eq:x-score}.}
From $x_t = \alpha_t x_0 + \sigma_t\varepsilon$:
\[
x^* = \E[x_0\mid x_t] = \frac{1}{\alpha_t}\big(x_t - \sigma_t\,\E[\varepsilon\mid x_t]\big) = \frac{1}{\alpha_t}\big(x_t + \sigma_t^2\,s^*\big),
\]
using~\eqref{eq:eps-score} to substitute $\varepsilon^* = -\sigma_t s^*$.

\emph{Equation~\eqref{eq:v-xeps}.}
The conditional velocity is $\dot{x}_t = \alpha_t' x_0 + \sigma_t'\varepsilon$.
Taking the expectation given $x_t$ yields $v^* = \alpha_t' x^* + \sigma_t'\varepsilon^*$ by linearity.
\end{proof}

\begin{remark}[Connection to posterior means]
Equations~\eqref{eq:eps-score}--\eqref{eq:x-score} are equivalent to: $\E[X_0\mid \bar{X}_t=x] = -\sigma_t\,s(x,t)$ and $\E[Z\mid\bar{X}_t=x] = \frac{1}{\alpha_t}(x+\sigma_t^2\,s(x,t))$ in the notation of Assumption~\ref{ass:affine-gauss} (where $X_0$ is noise and $Z$ is data).
\end{remark}

Their training objectives all share the template:
\begin{equation}\label{eq:general-loss}
\mathcal{L}(\theta) = \E_{x_0,\,\varepsilon,\,t}\!\Big[\omega(t)\,\big\|\mathrm{NN}_\theta(x_t,t) - (A_t x_0 + B_t\varepsilon)\big\|^2\Big].
\end{equation}

\begin{table}[h]
\centering
\renewcommand{\arraystretch}{1.3}
\begin{tabular}{@{} l c c c @{}}
\toprule
\textbf{Prediction type} & $A_t$ & $B_t$ & \textbf{Regression target} \\
\midrule
Clean ($x$-pred.) & $1$ & $0$ & $x_0$ \\
Noise ($\varepsilon$-pred.) & $0$ & $1$ & $\varepsilon$ \\
Score & $0$ & $-1/\sigma_t$ & $-\varepsilon/\sigma_t$ \\
Velocity ($v$-pred.) & $\alpha_t'$ & $\sigma_t'$ & $\alpha_t' x_0 + \sigma_t'\varepsilon$ \\
\bottomrule
\end{tabular}
\caption{Regression target $= A_t x_0 + B_t\varepsilon$ for each parameterization.}
\label{tab:four-param}
\end{table}

\subsection{The FM velocity is the PF-ODE drift: the complete equivalence}\label{subsec:fm-eq-diff}

We now derive the central identity connecting FM and diffusion.

\begin{theorem}[FM velocity = PF-ODE drift]\label{thm:fm-eq-pf}
Under Assumption~\ref{ass:affine-gauss}, the FM optimal velocity $v^*(x,t) = \E[\dot{\bar{X}}_t\mid \bar{X}_t = x]$ satisfies
\begin{equation}\label{eq:vel-score}
v^*(x,t) = \frac{\dot\alpha_t}{\alpha_t}\,x + \sigma_t\!\Big(\frac{\dot\alpha_t}{\alpha_t}\sigma_t - \dot\sigma_t\Big)s(x,t).
\end{equation}
Using the schedule conversion $f_t = \dot\alpha_t/\alpha_t$ and $g_t^2 = 2\sigma_t\dot\sigma_t - 2f_t\sigma_t^2$ (Proposition~\ref{prop:inverse-id}), this is exactly:
\begin{equation}\label{eq:v-eq-pf}
\boxed{v^*(x,t) = f_t\,x - \tfrac{1}{2}g_t^2\,\grad\log p_t(x),}
\end{equation}
which is the probability-flow ODE drift from Theorem~\ref{thm:pf-ode}.
\end{theorem}

\begin{proof}
From~\eqref{eq:affine-path}, $\dot{\bar{X}}_t = \dot\sigma_t X_0 + \dot\alpha_t Z$.
Conditioning on $\bar{X}_t = x$ and using the posterior mean identities from Proposition~\ref{prop:four-equiv}:
\begin{align*}
v^*(x,t) &= \dot\sigma_t\,\E[X_0\mid x] + \dot\alpha_t\,\E[Z\mid x] \\
&= \dot\sigma_t\,(-\sigma_t\,s) + \dot\alpha_t\cdot\frac{1}{\alpha_t}(x + \sigma_t^2\,s) \\
&= \frac{\dot\alpha_t}{\alpha_t}\,x + \Big(\frac{\dot\alpha_t}{\alpha_t}\sigma_t^2 - \dot\sigma_t\sigma_t\Big)s,
\end{align*}
which is~\eqref{eq:vel-score}.
Now substitute $f_t = \dot\alpha_t/\alpha_t$:
\[
v^*(x,t) = f_t\,x + \big(f_t\sigma_t^2 - \dot\sigma_t\sigma_t\big)\,s = f_t\,x - \frac{1}{2}\big(2\dot\sigma_t\sigma_t - 2f_t\sigma_t^2\big)\,s = f_t\,x - \frac{g_t^2}{2}\,s.
\]
\end{proof}

\begin{resultbox}{Gaussian FM = Diffusion: what this theorem means}
Under affine Gaussian conditional paths:
\begin{enumerate}[nosep]
\item The FM optimal velocity field \emph{is} the probability-flow ODE drift --- the same ODE used for deterministic sampling in diffusion models.
\item Training FM with $v$-prediction (CFM loss with target $\alpha_t' x_0 + \sigma_t'\varepsilon$) is gradient-equivalent to training diffusion with score/noise prediction, up to a time-dependent reweighting (since the four predictions are algebraically interchangeable by Proposition~\ref{prop:four-equiv}).
\item FM and diffusion are not two competing frameworks: they are the \emph{same mathematical object} viewed through different lenses (conditional velocity fields vs.\ score functions).
\end{enumerate}
\end{resultbox}

\subsection{PF-ODE in all four parameterizations}

Using the inter-conversion formulas from Proposition~\ref{prop:four-equiv}, the PF-ODE $\dot{x} = v^*(x,t)$ admits equivalent forms depending on which network output is available.

\begin{proposition}[PF-ODE forms]\label{prop:pf-ode-param}
\begin{align}
\frac{\dd x}{\dd t} &= \frac{\alpha_t'}{\alpha_t}x - \sigma_t\Big(\frac{\alpha_t'}{\alpha_t} - \frac{\sigma_t'}{\sigma_t}\Big)\varepsilon^*(x,t) & &(\varepsilon\text{-prediction}) \label{eq:pf-eps}\\
&= \frac{\sigma_t'}{\sigma_t}x + \alpha_t\Big(\frac{\alpha_t'}{\alpha_t} - \frac{\sigma_t'}{\sigma_t}\Big)x^*(x,t) & &(x\text{-prediction}) \label{eq:pf-x}\\
&= \frac{\alpha_t'}{\alpha_t}x + \sigma_t^2\Big(\frac{\alpha_t'}{\alpha_t} - \frac{\sigma_t'}{\sigma_t}\Big)s^*(x,t) & &(\text{score prediction}) \label{eq:pf-s}\\
&= v^*(x,t) & &(v\text{-prediction}). \label{eq:pf-v}
\end{align}
\end{proposition}

\begin{keybox}{The simplest parameterization}
In $v$-prediction, the PF-ODE is simply $\dot{x} = v^*(x,t)$: the velocity field \emph{is} the ODE drift.
This is why flow matching naturally trains a $v$-predictor, while diffusion traditionally trains an $\varepsilon$- or score-predictor and reconstructs the ODE drift via~\eqref{eq:pf-eps} or~\eqref{eq:pf-s}.
\end{keybox}

\begin{example}[VP simplification]
Under VP schedules $\alpha_t = \sqrt{\bar\alpha_t}$, $\sigma_t = \sqrt{1-\bar\alpha_t}$, define $\kappa_t = \dot\alpha_t/\alpha_t = \dot{\bar\alpha}_t/(2\bar\alpha_t)$.
Then the coefficient of $s$ in~\eqref{eq:vel-score} equals $\kappa_t$ as well:
$\sigma_t(\kappa_t\sigma_t - \dot\sigma_t) = \kappa_t$.
The FM velocity simplifies to $v^*(x,t) = \kappa_t\,(x + s(x,t))$.
\end{example}

% =====================================================================================
\section{All affine flows are equivalent}\label{sec:all-affine-equiv}

\begin{proposition}[Equivalence of affine flows]\label{prop:affine-equiv}
Let $x_t^{\mathrm{FM}} = (1-t)x_0 + t\varepsilon$ be the canonical FM interpolation and $x_t = \alpha_t x_0 + \sigma_t\varepsilon$ any affine interpolation.
Define $c(t) \coloneqq \alpha_t + \sigma_t$ and $\tau(t) \coloneqq \sigma_t/(\alpha_t + \sigma_t)$ (assuming $c(t)\neq 0$).
Then:
\begin{equation}\label{eq:affine-equiv}
x_t = c(t)\,x_{\tau(t)}^{\mathrm{FM}},
\end{equation}
and the velocity transforms as:
\begin{equation}\label{eq:affine-vel-equiv}
v(x_t,t) = c'(t)\,x_{\tau(t)}^{\mathrm{FM}} + c(t)\,\tau'(t)\,v^{\mathrm{FM}}\!\big(x_{\tau(t)}^{\mathrm{FM}},\tau(t)\big).
\end{equation}
\end{proposition}

\begin{proof}
Direct algebraic rewrite:
$x_t = (\alpha_t + \sigma_t)\big(\frac{\alpha_t}{\alpha_t+\sigma_t}x_0 + \frac{\sigma_t}{\alpha_t+\sigma_t}\varepsilon\big) = c(t)\big((1-\tau(t))x_0 + \tau(t)\varepsilon\big) = c(t)\,x_{\tau(t)}^{\mathrm{FM}}$.

For the velocity, apply the chain rule to $x_t = c(t)\,x_{\tau(t)}^{\mathrm{FM}}$ and use $v^{\mathrm{FM}}(x,\tau) = \varepsilon - x_0$ (constant along the canonical path).
\end{proof}

\begin{resultbox}{Practical consequence}
The ``simplicity'' of the linear FM schedule $(\alpha_t = 1-t,\,\sigma_t = t)$ is not fundamental: any affine interpolation is equivalent up to time reparameterization and spatial rescaling.
Different schedules affect training dynamics, weighting, and numerical properties, but not the class of representable flows.
\end{resultbox}

\newpage
% =====================================================================================
%  PART V: SYNTHESIS
% =====================================================================================

\part{Synthesis}

% =====================================================================================
\section{The Fokker--Planck equation as the unifying principle}\label{sec:synthesis}

We conclude by summarizing the logical architecture of the course in a single diagram.

\begin{summarybox}{Unified perspective}
All generative models in this course share one underlying rule: \textbf{the Fokker--Planck / continuity equation} governing the evolution of $p_t(x)$ under a chosen forward process.

\medskip
\textbf{Forward process} (from data to noise):
\begin{itemize}[nosep]
\item Variational: defining $p_t(x_t\mid x_0) = \cN(\alpha_t x_0, \sigma_t^2\Id)$.
\item SDE: $\dd x = f(t)x\dd t + g(t)\dd w$, with $(f,g) \leftrightarrow (\alpha_t,\sigma_t)$ via Proposition~\ref{prop:inverse-id}.
\end{itemize}

\medskip
\textbf{Reverse process} (from noise to data), all sharing the same $p_t$:
\begin{itemize}[nosep]
\item PF-ODE: $\dd x = v^*(x,t)\dd t$ where $v^* = f(t)x - \frac{1}{2}g^2(t)\grad\log p_t(x)$ \hfill ($\gamma\equiv 0$)
\item Reverse SDE: drift corrected by $-g^2 s$ plus diffusion $g\dd\bar{W}$ \hfill ($\gamma = g$)
\item General SDE: $\dd x = [v^* + \frac{\gamma^2}{2}s]\dd t + \gamma\dd B$ for any $\gamma(t)\ge 0$ \hfill (Thm.~\ref{thm:marginal-family})
\end{itemize}

\medskip
\textbf{Learning}:
\begin{itemize}[nosep]
\item Score matching $\leftrightarrow$ Denoising score matching $\leftrightarrow$ Conditional flow matching.
\item All are conditional expectations: $s^* = \E[\grad\log p_t(x_t\mid x_0)\mid x_t]$, $v^* = \E[\dot{x}_t\mid x_t]$.
\item Training with conditional targets differs from marginal targets by a $\theta$-free constant.
\end{itemize}

\medskip
\textbf{Equivalences} (Theorem~\ref{thm:fm-eq-pf} and Propositions~\ref{prop:four-equiv},~\ref{prop:affine-equiv}):
\begin{itemize}[nosep]
\item Four parameterizations ($\varepsilon$, $x$, score, velocity) are algebraically interchangeable.
\item All affine flows are equivalent up to time reparameterization and rescaling.
\item Gaussian FM $=$ Diffusion: the FM optimal velocity \emph{is} the PF-ODE drift $f_t x - \frac{1}{2}g_t^2\grad\log p_t$.
\end{itemize}
\end{summarybox}

\newpage
% =====================================================================================
%  APPENDICES
% =====================================================================================

\appendix

\section{Derivation of weak forms from stochastic calculus}\label{app:weak-forms}

\subsection{ODE case}
If $\dot{X}_t = v_t(X_t)$ and $\varphi\in C_c^\infty$, the chain rule gives $\frac{\dd}{\dd t}\varphi(X_t) = \inner{\grad\varphi(X_t)}{v_t(X_t)}$.
Taking expectations yields the weak continuity equation~\eqref{eq:cont-weak}.

\subsection{SDE case}
If $\dd X_t = b(X_t,t)\dd t + \sigma(t)\dd B_t$, It\^o's formula gives
\[
\dd\varphi(X_t) = \inner{\grad\varphi(X_t)}{b(X_t,t)}\dd t + \sigma(t)\inner{\grad\varphi(X_t)}{\dd B_t} + \frac{\sigma(t)^2}{2}\Delta\varphi(X_t)\dd t.
\]
Taking expectations (the It\^o integral contributes zero) yields the weak Fokker--Planck equation~\eqref{eq:fp-weak}.

\end{document}
